\section{Fakultät}
\label{sec:fundamentals_factorial}

Die Fakultät ist insbesondere beim Zählen von Anordnungen und
Auswahlmöglichkeiten von Bedeutung und bildet später die Grundlage für die
Binomialkoeffizienten in Abschnitt \ref{subsec:fundamentals_general_binomial_theorem}.

\begin{DefinitionBox}{Fakultät}{fundamentals_factorial}
  Für $n\in\N$ ist die \emph{Fakultät} \EnglishTerm{factorial} von $n$
  definiert durch
  \[
    n!=\prod_{k=1}^{n}k=1\cdot2\cdot3\cdots n.
  \]
  Zusätzlich gilt: $0!=1.$
\end{DefinitionBox}

\begin{ExampleBox}{Erste Fakultäten}{fundamentals_first_factorials}
  \begin{enumerate}[label=\alph*),leftmargin=2em,itemsep=0.8em]
    \item Fakultäten von $0!$ bis $5!$:
      \begin{ClozeArea}[28]{12}
        \[
          \begin{array}{r@{\;}c@{\;}l}
            0!&=&1,\\[0.5em]
            1!&=&1,\\[0.5em]
            2!&=&1\cdot2=2,\\[0.5em]
            3!&=&1\cdot2\cdot3=6,\\[0.5em]
            4!&=&1\cdot2\cdot3\cdot4=24,\\[0.5em]
            5!&=&1\cdot2\cdot3\cdot4\cdot5=120.
          \end{array}
        \]
      \end{ClozeArea}

    \item Quotient zweier konkreter Fakultäten:
      \begin{ClozeArea}[28]{4}
        \[
          \frac{7!}{6!}
            =\frac{1\cdot2\cdots6\cdot7}{1\cdot2\cdots6}
            =7.
        \]
      \end{ClozeArea}

    \item Allgemeiner Quotient aufeinanderfolgender Fakultäten:
      \begin{ClozeArea}[28]{4}
        \[
          \frac{(n+1)!}{n!}
            =\frac{1\cdot2\cdots n\cdot(n+1)}{1\cdot2\cdots n}
            =n+1.
        \]
      \end{ClozeArea}
  \end{enumerate}
\end{ExampleBox}

\begin{KeyPointBox}{Rekursive Berechnung}{fundamentals_factorial_recursion}
  Für $n\in\N_0$ gilt
  \[
    (n+1)!=\Cloze{8cm}{(n+1)\,n!}.
  \]
  Damit lässt sich jede Fakultät aus der vorhergehenden berechnen.
\end{KeyPointBox}
